6. Discussion
6.1 Full-Year Calibration Masks Event-Window Undercoverage
The full-year 2014 QR-GBT results present a broadly acceptable calibration picture: the nominal 90% prediction interval achieves an empirical coverage rate of 86.8%, and the nominal 98% prediction interval achieves 97.3%. Evaluated solely on annual aggregate metrics, the probabilistic forecasts appear reasonably well-calibrated. Disaggregation by evaluation window, however, reveals that this aggregate summary conceals substantial undercoverage during the most operationally consequential period in the evaluation year.
During the January 6–8 vortex window, the nominal 90% PI empirical coverage falls to 66.7% — a decline of 20.1 percentage points from the full-year rate and 23.3 percentage points below the nominal 90% target. The nominal 98% PI coverage falls to 84.7%, more than 12 percentage points below its full-year rate. The top 1% of load hours by observed load level shows consistent undercoverage: 90% PI coverage of 63.2% and 98% PI coverage of 82.8%. Coverage deteriorates during the conditions where accurate uncertainty quantification would be most consequential.
This pattern is consistent with a general risk in aggregate calibration assessment: when a high-frequency, moderate-load regime dominates the annual sample, aggregate coverage rates can appear adequate while masking systematic failures in lower-frequency, high-load regimes. In this case study, the full-year 2014 sample contains 8,760 hours; the vortex window represents only 72 of those hours, and the top 1% subset approximately 88. The contribution of these hours to the aggregate coverage statistic is small, yet they represent the tail-risk conditions that probabilistic forecasts are most expected to address. These findings suggest that evaluation designs relying solely on full-year aggregate metrics may not surface stress-window calibration failures, and that window-specific coverage reporting provides material additional information.
6.2 Extreme Cold as a Distribution-Shift Stress Test
The January 6–8 polar vortex event constitutes a near-annual-peak cold-weather stress event within the 2014 evaluation year. The associated load conditions — high sustained demand driven by intense cold, combined with near-peak load magnitudes — are consistent with distribution-shift stress relative to more moderate operating conditions that make up the more moderate operating conditions that are expected to dominate many multi-year hourly samples. Whether the 2010–2013 training period contains weather events of comparable intensity is not directly verified here; this is addressed as a formal limitation in Section 7. The observed degradation in model performance during this window is consistent with the cold-event regime being underrepresented relative to more moderate operating conditions in the training data, but direct training-distribution support for this interpretation is not verified by this study.
All weather-informed model results use same-hour ERA5 retrospective reanalysis weather input. The performance degradation is observable across point and probabilistic forecast approaches: GBoost point forecast MAE increases from 721 MW full-year to 1,705 MW during the vortex window, approximately 2.4 times the full-year rate. QR-GBT q50 MAE increases from 761 MW to 2,130 MW. Mean pinball loss increases from 153 to 473.
The relative degradation of Naive-24h and Naive-168h is more severe, with vortex-window MAE reaching 15,598 MW and 24,531 MW respectively. These baselines anchor to load observed under different thermal conditions one day and one week prior, making them structurally ill-suited to events in which the driving weather conditions depart sharply from the recent past. Linear Regression, despite using the same retrospective ERA5 weather inputs as GBoost, produces a vortex MAE of 4,248 MW versus 1,705 MW for GBoost, indicating that the linear specification was less able than GBoost to represent event-window behavior under extreme cold conditions by a linear model in this case study.
6.3 Why Point Accuracy and Probabilistic Coverage Diverge
A notable pattern in the results is that GBoost point forecast accuracy, while degraded during the vortex window, remains substantially lower in absolute error than the naive and linear baselines. GBoost vortex MAE is 1,705 MW, compared with 4,248 MW for Linear Regression and 15,598 MW for Naive-24h. Yet at the same time, the QR-GBT probabilistic model — which uses the same feature set and training data — produces 90% PI coverage of only 66.7% during the vortex window and fails to bound the event peak even within the nominal 98% prediction interval.
This divergence illustrates that point forecast accuracy and probabilistic interval calibration are distinct properties of a forecasting system, and that improvements in one do not necessarily imply improvements in the other. A model can achieve relatively low median error while still failing to assign sufficient probability mass to the tails of the conditional load distribution.
At the event peak hour (Jan 7 18:00 EPT), the observed load exceeded the post-rearrangement QR-GBT q99 estimate of 139,585 MW, leaving the event peak outside the nominal 98% prediction interval. The q99 estimate remained 925 MW below the observed peak of 140,510 MW. The q50 estimate at this hour was 135,357 MW, placing it 5,153 MW below the observed peak. These figures illustrate that the upper quantile estimates at the event peak hour were insufficient to bound the observed outcome, independent of the median forecast accuracy.
It is worth noting that this finding pertains specifically to the behavior of independently fitted quantile regression models evaluated against a near-annual-peak cold-weather stress event in this case study. It does not generalize as a statement about quantile regression methods broadly, and does not imply that the QR-GBT model is without value across the full evaluation year.
6.4 Implications of the Winter–Summer Near-Peak Comparison
The winter–summer comparison is designed to distinguish between two possible explanations for event-window performance degradation: general peak-load difficulty and cold-event-specific difficulty. The two windows have similar peak-load magnitudes — 140,510 MW for the winter vortex and 141,678 MW for the summer near-peak window, a difference of less than 1% of the summer peak — allowing load magnitude to be approximately controlled in the comparison.
Despite this similarity, probabilistic performance is materially worse during the cold-event window across all reported metrics. The q50 MAE during the vortex window (2,130 MW) is approximately twice that of the summer window (1,060 MW). Mean pinball loss is 473 during the vortex window versus 209 during the summer window. The Winkler score for the nominal 90% interval is 15,476 MW during the vortex window versus 6,355 MW during the summer window. Nominal 90% PI coverage is 66.7% during the vortex window and 86.1% during the summer window.
The summer window 90% PI coverage (86.1%) is close to the full-year rate (86.8%), suggesting that summer near-peak conditions do not present the same forecasting challenge as the winter cold-event window, at least in this case study. The comparison suggests that cold-event conditions posed a greater challenge than the summer near-peak window in this case study. A direct explanation in terms of training-distribution support — such as whether summer peak-load regimes are better represented in the training data than winter cold-event regimes — would require additional analysis of training-period weather and load distributions and should be treated as a limitation rather than a conclusion of this study.
6.5 Quantile Crossing and Post-Hoc Rearrangement
Quantile crossing was detected in 5,786 of 8,760 full-year 2014 test hours, representing approximately 66% of all test hours. This high crossing rate indicates that the independently fitted QR-GBT quantile models — each optimizing the pinball loss at a single quantile level without joint monotonicity constraints — produce inconsistent orderings across a substantial fraction of the evaluation period. Post-hoc monotonic rearrangement was applied to restore a valid ordering, and all reported prediction intervals use rearranged estimates. Rearrangement resolves the ordering violation as a technical matter but does not address the underlying cause of the crossing and does not guarantee well-calibrated intervals.
The crossing rate of 66% suggests that the seven independently fitted models do not maintain coherent monotonic ordering in the majority of test hours. This is a structural limitation of the independently fitted quantile approach and is consistent with findings in the broader quantile regression literature [CITATION: quantile crossing reference]. In applications where a coherent predictive distribution is required, post-hoc rearrangement may be insufficient as a standalone correction, and joint or distributional estimation approaches may be more appropriate.
The reported event-peak q95 and q99 values are post-rearrangement values; therefore, interpretation should focus on the fact that the event peak remains outside the rearranged nominal 98% interval, rather than on unverified claims about the original raw quantile ordering at that hour. The verified post-rearrangement values — q95 of 139,410 MW and q99 of 139,585 MW against an observed peak of 140,510 MW — are the basis for the tail-risk conclusions reported in Section 5.6, and the discussion here does not extend beyond those verified figures.
6.6 Retrospective Weather Inputs and Benchmark Interpretation
All GBoost and QR-GBT model results in this study use same-hour ERA5 retrospective reanalysis weather input. ERA5 reanalysis provides high-quality gridded weather estimates derived from a consistent data assimilation system applied retrospectively to historical observations, but it is not equivalent to an operational numerical weather prediction forecast issued in advance of the target hour. Same-hour ERA5 inputs are not available in real time and represent an idealized information setting that does not reflect the uncertainty introduced by weather forecast errors under operational conditions.
PJM Day-Ahead forecasts are produced operationally using PJM's internal forecasting system and inputs not controlled by this study. Comparisons between GBoost and PJM Day-Ahead — including the full-year MAE comparison (721 MW versus 1,937 MW) and the vortex-window MAE comparison (1,705 MW versus 3,148 MW) — are therefore retrospective in character. The gradient boosted models in this study have access to same-hour retrospective reanalysis weather information that was not available at the time of operational forecast issuance at the time of forecast issuance. These comparisons provide useful retrospective context for interpreting the magnitude of the modeled errors relative to those of an operational forecasting system during the same period, but they do not constitute evidence of operational superiority.
The informational asymmetry between retrospective ERA5 inputs and operational forecast-weather inputs is a fundamental constraint on the generalizability of the results reported in Section 5. Quantifying the degradation in model performance that would arise from replacing retrospective ERA5 inputs with real-time NWP forecast-weather inputs is outside the scope of this study. Any extension of these results toward an operational context would require separate validation using forecast-weather inputs and a prospective evaluation design; that scope is deferred to Section 7 as a formal limitation.